\section{Conclusion}\label{sec:conclusion}

Whether a machine learning model is ``better'' than another is a statistical claim, and on spatio-temporal data it depends on two choices that are rarely made explicit: the validation design, which defines what is estimated, and the inferential procedure, which defines how its uncertainty is assessed. We formalised the comparison of two models as inference on the mean of a dependent loss-differential field with correlated fold effects, showed why the paired $t$-test on fold scores fails under random cross-validation, and proposed DA-HAC, a test that combines a variogram-calibrated product-kernel variance, a training-variance term and a Satterthwaite reference distribution, all computed from one blocked cross-validation run.

In controlled simulations and with gradient-boosting learners, conventional procedures declared false improvements up to sixteen times more often than their nominal level, and random cross-validation could halve the estimated gain of a genuinely better model. Blocked cross-validation combined with DA-HAC, or with the paired $t$-test on blocked folds, brought the error rate close to the nominal level under weak dependence and reduced it substantially under stronger dependence, without fully restoring it when the dependence range approaches the size of the study domain, a regime that DA-HAC's diagnostics help to detect. On three public data sets, the conclusion of four of the six comparisons depended on the validation design or on the test: an apparent superiority of XGBoost on hourly weather data vanished under temporal blocking, and two smaller gains were declared significant by independence-based procedures but not by DA-HAC, whose diagnostics showed that one of them could not be resolved with the available data. Large improvements, such as those of gradient boosting over linear baselines, were confirmed by every procedure.

Trustworthy machine learning requires not only accurate and explainable models but also trustworthy evidence that one model is better than another. We hope that the procedure, the diagnostics and the open-source implementation presented here will help practitioners in environmental, agricultural, epidemiological and mobility applications to report improvements that are statistically real.
